\subsection{"捡尸"与酒后失能：旁观者能做什么}

"捡尸"是中文语境里对\textbf{趁他人严重醉酒丧失意识之际实施性行为}的俗称。它不是一个边缘现象——它发生在酒吧、夜店、KTV 与集会散场后的街边，而且\textbf{绝大多数情况下本可以被身边的人阻止}。这一节写给所有人：不只是"提醒你别喝醉"，而是"告诉你旁观者能做什么"。

\textbf{一、先把法律讲清楚}
\begin{itemize}
\item 酒精不构成免责事由。与丧失意识或无法表达意愿的人发生性行为，属\textbf{强奸}，无论对方是否"先前愿意"、是否"没有反抗"；
\item 判断标准只有一个：\textbf{此刻对方是否具备表达同意的能力}。醉到说不了完整句子、走不了直线、喊不醒——这种情况下不存在"同意"；
\item 常见辩解"她跟我走的""她没拒绝"在法律上不成立。"跟我走"不等于"同意性行为"，"没拒绝"不等于"同意"——这两点在近年司法实践中已有明确裁判倾向。
\end{itemize}

\textbf{二、旁观者干预：5D 策略}
\textbf{（一）Direct（直接介入）}：最有效也最需要勇气的方式。走上前对当事人说"你朋友在找你""我送你回去"，用一句自然的话切断正在发生的事——不需要对抗，只需要打断；
\textbf{（二）Distract（转移注意力）}：假装问路、假装认错人、把一个"聚会理由"抛给施害者，为当事人争取离开的时间；
\textbf{（三）Delegate（求助）}：找场所保安、店员、同伴，或直接报警。场所经营者有义务保障场内安全，"怕多事"的风险远小于放过的后果；
\textbf{（四）Delay（拖延）}：在当事人被带走前尽可能拖住——搭话、提出一起走、故意制造一点麻烦。多数情形下，施害者最怕的就是"有人盯着"；
\textbf{（五）Document（记录）}：在安全的前提下，记下时间、地点、人物、车牌、衣着特征，拍下（不传播）能帮助识别施害者的信息。记录的目的是维权，不是围观。

\textbf{三、同伴之间可执行的约定}
\begin{itemize}
\item \textbf{结伴原则}：出门前约定"不单独离场"，离场必须两人以上同行并互相确认；\textbf{手中饮品不离视线}，换杯后重新确认；
\item \textbf{互相盯防}：约定一个"我不对劲"的暗号（一句话、一个表情），收到即立刻带走，事后不追问、不埋怨、不传播——不被指责的环境，是朋友愿意求助的前提；
\item \textbf{送人到家}：交到可信的人手上或确认进屋关门再离开，中途不"顺路交给别人"；
\item \textbf{场所与交通}：优先选择有正规安保与监控的场所；代驾、网约车的行程分享给同伴；深夜打车坐后排并核对车牌。
\end{itemize}

\textbf{四、若事情已经发生}
\begin{itemize}
\item \textbf{先做医学处置，再做其他决定}：72 小时内评估 HIV 阻断（PEP）、紧急避孕与性传播感染预防；保留衣物、不洗澡、不冲洗，尽快就医取证；
\item 报警与否由当事人决定，但取证与医学处理不建议等待——错过时间窗之后，可用的证据会大幅减少；
\item 对当事人：\textbf{不指责"为什么喝那么多"、不追问细节、不传播}。责任永远在施害者一方。
\end{itemize}

\noindent\textit{【编者按】}"捡尸"一词本身带着对受害者的戏谑，把严重的性侵包装成一个猎奇笑谈。使用这个词时值得想一想：它描述的是一起可能判处有期徒刑的犯罪。
